% Industry resume, Grover Aruquipa
% Compile with pdfLaTeX (also works on Overleaf). One column on purpose:
% applicant tracking systems read single-column text far more reliably than sidebars.
\documentclass[10pt,a4paper]{article}
\usepackage[T1]{fontenc}
\usepackage[utf8]{inputenc}
\usepackage[scaled=0.95]{helvet}
\renewcommand{\familydefault}{\sfdefault}
\usepackage[top=1.1cm,bottom=1.0cm,left=1.5cm,right=1.5cm]{geometry}
\usepackage{enumitem}
\usepackage{titlesec}
\usepackage{xcolor}
\usepackage{microtype}
\usepackage[hidelinks]{hyperref}

\definecolor{ink}{HTML}{111111}
\definecolor{accent}{HTML}{0D2B61}
\definecolor{muted}{HTML}{555555}
\color{ink}
\pagestyle{empty}
\setlength{\parindent}{0pt}
\setlength{\parskip}{0pt}

\titleformat{\section}{\normalsize\bfseries\color{accent}}{}{0em}{\MakeUppercase}[\vspace{-6pt}{\color{accent}\rule{\linewidth}{0.6pt}}]
\titlespacing*{\section}{0pt}{7pt}{3pt}
\setlist[itemize]{leftmargin=1.1em,itemsep=0.5pt,topsep=1pt,parsep=0pt,label=\textbullet}

% \role{Title}{Organization}{Location}{Dates}
\newcommand{\role}[4]{\textbf{#1}, #2 \hfill {\color{muted}#3 \textbar\ #4}\par}

\begin{document}

%================ HEADER ================
{\LARGE\bfseries Grover Aruquipa}\par\vspace{2pt}
{\large R\&D Mechanical Engineer \textbar\ Precision Robotic Joints and Wrists, Concept to Verified Prototype}\par\vspace{3pt}
{\small Qu\'ebec, QC, Canada \ \textbullet\ \
\href{mailto:grover.aruquipa.9@gmail.com}{grover.aruquipa.9@gmail.com} \ \textbullet\ \
\href{https://www.linkedin.com/in/grover-aruquipa-996917126}{linkedin.com/in/grover-aruquipa-996917126} \ \textbullet\ \
\href{https://groveraruquipa.github.io}{groveraruquipa.github.io}}\par


%================ SUMMARY ================
\section{Summary}
Mechanical engineer finishing a PhD in robot mechanism design (Universit\'e Laval, defense early 2027). I have taken three
precision robots from concept and SolidWorks design through joint and transmission design, component selection,
fabrication, assembly and test to validated hardware, and I am building a fourth. The most compact is a 4-DoF coaxial
wrist-gripper that rolls its tool without limit while the jaw is driven through concentric hollow shafts, the same
problem a wristed surgical instrument has to solve. Working in English and French; ready to relocate to Vaud.

%================ SKILLS ================
\section{Technical Skills}
\begin{itemize}
\item \textbf{Mechanical design:} SolidWorks (daily use for three full robot assemblies), Creo; precision revolute and
prismatic joints, coaxial and concentric transmissions (hollow shafts, belts, gears), bearing arrangements, integrated
grippers; motor, encoder, gearbox and bearing selection from vendor catalogs.
\item \textbf{Analysis:} FEA in ANSYS; rigid-body kinematics and dynamics, Jacobians, singularity and workspace analysis;
simulation in MATLAB, Simscape and MuJoCo.
\item \textbf{Prototyping and verification:} design for machining and for SLS and FDM printing; machining, printing,
assembly, wiring and commissioning of prototypes; test setups for tracking accuracy and repeatability; real-time
multi-axis control at 500\,Hz (OPAL-RT, Simulink Real-Time) to run the tests.
\item \textbf{Software and electronics:} MATLAB, Python, C++ (real time), ROS~2, Pinocchio; Altium (basic PCB),
motor and encoder wiring.
\end{itemize}

%================ EXPERIENCE ================
\section{Experience}
\role{Doctoral Researcher}{Robotics Laboratory, Universit\'e Laval}{Qu\'ebec, Canada}{2023 -- present}
\begin{itemize}
\item Owned the full mechanical development of three multi-DoF robots (4, 7 and 8 actuators, all base-mounted, gripper
driven through the mechanism): requirements, concept, SolidWorks design, component selection, fabrication, assembly,
test and redesign.
\item \textbf{4-DoF coaxial wrist-gripper:} four actuation channels through concentric hollow shafts, non-backdrivable
gripper, 180\,mm footprint, 0.25\,kg moving mass; unlimited tool-axis roll verified over ten continuous revolutions with
sub-degree tracking under 500\,Hz real-time control.
\item \textbf{(6+1)-DoF parallel robot:} seven base-mounted actuators; gripper actuated from the base through a
concentric tube while the platform rotates without limit.
\item \textbf{Four-leg Delta robot with integrated Sarrus gripper (in progress):} grasp force estimated from motor
currents, without a force sensor at the jaw, to pick fragile objects quickly.
\item Presented every design to supervisors and co-authors before fabrication; day-to-day lead of two graduate students
on an 8-actuator robot. First author of 3 journal papers, 4 under review; reviewer for ASME J.\ Mech.\ Robotics.
\end{itemize}

\vspace{3pt}
\role{Visiting Researcher}{Institute of Robotics, Johannes Kepler University Linz}{Linz, Austria}{Apr. -- May 2026}
\begin{itemize}
\item Developed singularity and grasp-quality indices for manipulators with integrated grasping, hosted by Univ.-Prof.~A.~M\"uller (Mitacs Globalink Research Award).
\end{itemize}

\vspace{3pt}
\role{Research Intern}{FEMTO-ST (CNRS)}{Besan\c{c}on, France}{2022 -- 2023}
\begin{itemize}
\item Designed and modeled a redundant dexterous parallel robot with two unlimited rotations for micro-manipulation, and wrote its real-time C++ control software; clean-room fabrication projects.
\end{itemize}

\vspace{3pt}
\textbf{Industry roles} \hfill {\color{muted}2020 -- 2021}\par
{\small Computer Vision Engineer, Sencon Srl (industrial vision, scaffolding detection) \ \textbullet\ \ Data Science
Intern, Hitachi ABB Power Grids, Zurich (machine learning for high-voltage circuit breakers) \ \textbullet\ \ Data
Scientist, KTH Royal Institute of Technology.}

%================ EDUCATION ================
\section{Education}
\textbf{PhD, Mechanical Engineering}, Universit\'e Laval \hfill {\color{muted}Qu\'ebec, Canada \textbar\ 2023 -- 2027 (expected)}\par
{\small Kinematically redundant parallel robots with unlimited rotations and integrated grasping. Advisors: C.~Gosselin, P.~Lambert.}\par\vspace{2pt}
\textbf{MSc, Smart Integrated Systems}, Universit\'e Bourgogne--Franche-Comt\'e \hfill {\color{muted}Besan\c{c}on, France \textbar\ 2021 -- 2023}\par\vspace{2pt}
\textbf{BEng, Mechatronics Engineering}, Bolivian Catholic University \hfill {\color{muted}La Paz, Bolivia \textbar\ 2014 -- 2019}\par
{\small Thesis: reconfigurable Delta robot for product picking.}

%================ AWARDS ================
\section{Awards and Languages}
\begin{itemize}
\item 1st Place, Graduate Division, ASME Student Mechanisms and Robotics Design Competition (2025), team lead.
\item Finalist, MassRobotics competition, Boston, with one of my robot prototypes, team lead.
\item Mitacs Globalink Research Award (2026); 1st Place, Student Paper Competition, Argentine Robotics Conference (2019).
\item Languages: Spanish (native), English (full professional), French (intermediate).
\end{itemize}

\newpage
%================ PUBLICATIONS ================
\section{Publications}
{\small First author on all papers.}\par\vspace{2pt}
\textbf{Journal articles}
\begin{itemize}
\item \textbf{G.~Aruquipa}, P.~Lambert, R.~Dahmouche, C.~Gosselin, ``Parallel robots for grasping tasks: a review,''
\emph{ASME Journal of Mechanisms and Robotics}, vol.~18, no.~12, art.~120801, 2026. doi:10.1115/1.4072612
\item \textbf{G.~Aruquipa}, V.~Muralidharan, P.~Lambert, C.~Gosselin, ``Analysis and design of a novel compact
four-degree-of-freedom parallel robot with unlimited rotation,'' \emph{ASME Journal of Mechanisms and Robotics},
vol.~18, no.~6, art.~061001, 2026.
\item \textbf{G.~Aruquipa}, P.~Lambert, C.~Gosselin, ``Analysis and design of a novel (6+1)-DoF parallel robot with
unlimited rotation and grasping capabilities,'' \emph{Mechanism and Machine Theory}, vol.~224, art.~106438, 2026.
\end{itemize}
\textbf{Conference papers}
\begin{itemize}
\item \textbf{G.~Aruquipa}, P.~Lambert, C.~Gosselin, ``Kinematic analysis and design of a novel 9-DoF parallel robot
with grasping capabilities,'' \emph{CCToMM Symposium on Mechanisms, Machines and Mechatronics}, 2025, pp.~117--127.
\item \textbf{G.~Aruquipa}, F.~Diaz, ``An IoT architecture based on the control of a bio-inspired manufacturing system
for the detection of anomalies with vibration sensors,'' \emph{Procedia Computer Science}, vol.~200, pp.~438--450, 2022.
\end{itemize}
\textbf{Under review}
\begin{itemize}
\item \textbf{G.~Aruquipa}, P.~Lambert, C.~Gosselin, ``A kinematically redundant (6+2)-DoF parallel robot with three
orthogonal unlimited rotations and coaxial actuation,'' submitted to \emph{IEEE ICRA 2027}.
\item \textbf{G.~Aruquipa}, P.~Lambert, C.~Gosselin, ``A compact coaxial parallel wrist-gripper with unlimited
tool-axis rotation,'' submitted to \emph{Robotics and Autonomous Systems}.
\item \textbf{G.~Aruquipa}, P.~Lambert, C.~Gosselin, ``An analytical framework for collision-free and singularity-free
unlimited rotation in spatial parallel mechanisms,'' submitted to \emph{Mechanism and Machine Theory}.
\item \textbf{G.~Aruquipa}, A.~Ayoub, R.~Dahmouche, ``Architecture, kinematics, and workspace analysis of a redundant
8RSS parallel robot with two unlimited rotations,'' in revision, \emph{Mechanism and Machine Theory}.
\end{itemize}


%================ REFERENCES ================
\section{References}
\begin{itemize}
\item \textbf{Prof.~Cl\'ement Gosselin}, Full Professor, Department of Mechanical Engineering, Universit\'e Laval, Qu\'ebec,
Canada. PhD supervisor. \href{mailto:clement.gosselin@gmc.ulaval.ca}{clement.gosselin@gmc.ulaval.ca}
\item \textbf{Prof.~Patrice Lambert}, Assistant Professor, Department of Mechanical Engineering, Universit\'e Laval,
Qu\'ebec, Canada. PhD co-supervisor. \href{mailto:patrice.lambert@gmc.ulaval.ca}{patrice.lambert@gmc.ulaval.ca}
\item \textbf{Univ.-Prof.~Dr.-Ing.~Andreas M\"uller}, Head of the Institute of Robotics, Johannes Kepler University Linz,
Austria. Host during my research stay (2026). \href{mailto:a.mueller@jku.at}{a.mueller@jku.at}
\end{itemize}

\end{document}
